% Seção: Metodologia  (orçamento: 2.5 página(s))
%
% Referências previstas (chaves de refs.bib): schroer2021crispdm, martinezplumed2021crispdm, shimaoka2024crispdm, decroos2019actions, brandes2001faster, yao2023react, ozsoy2025text2cypher, lewis2020rag, sainz2023nlp, ji2023survey
% Figuras e tabelas ainda a fazer: F1 (arquitetura), F5 (exemplo de robustez).

\section{Metodologia}\label{sec:metodologia}

O trabalho seguiu as seis fases do CRISP-DM~\cite{schroer2021crispdm}: entendimento do problema, entendimento
dos dados, preparação, modelagem, avaliação e implantação. A pergunta de partida é quem faz o cálculo quando um modelo de
linguagem responde sobre dados de eventos de futebol: o próprio modelo, uma consulta que ele escreve ou código que
já está escrito.

\subsection{Dados}

Os dados são os eventos abertos da StatsBomb~\cite{statsbomb_opendata} da final da Copa do Mundo de 2022. Os
\nEvents{} eventos da partida foram convertidos em \nActions{} ações no formato SPADL~\cite{decroos2019actions},
em uma tabela com uma linha por ação, com zona do campo, fase de posse e valor esperado de ameaça (xT). A tabela
foi carregada em um grafo Neo4j. Uma única partida foi usada para que a tabela inteira caiba no
contexto do modelo: ela ocupa \tokInEventsInPrompt{} tokens por pergunta.

\subsection{Formas de acesso}

Foram comparadas \nArms{} formas de acesso (braços), ordenadas por quanta lógica de consulta já está pronta antes
de o modelo agir (Tabela~\ref{tab:bracos}). Todas usam o mesmo modelo (\modelName{}, temperatura zero), o mesmo
prompt de sistema, que não nomeia a partida, e a mesma saída estruturada (jogadores, valor, indicação de ausência
de dados e justificativa).

\begin{table}[H]
\centering\small
\begin{tabular}{p{0.34\columnwidth}p{0.58\columnwidth}} % lint: ignore
\toprule
Braço & O que o modelo recebe \\
\midrule
\armNoContext{} (\texttt{no\_context}) & só a pergunta \\
\armVector{} (\texttt{vector}) & as \nTopK{} linhas de eventos recuperadas por similaridade de vetores \\
\armEventsInPrompt{} (\texttt{events\_in\_prompt}) & a tabela de \nActions{} ações inteira \\
\armTextToCypher{} (\texttt{text\_to\_cypher}) & o esquema do grafo e uma ferramenta que executa a consulta Cypher escrita por ele \\
\armGraphTools{} (\texttt{graph\_tools}) & \nTools{} ferramentas parametrizadas sobre o grafo \\
\armStatsInPrompt{} (\texttt{stats\_in\_prompt}) & totais por jogador e por time, já calculados \\
\bottomrule
\end{tabular}
\caption{As \nArms{} formas de acesso, na ordem da escada.}\label{tab:bracos}
\end{table}

O \armVector{} segue o padrão de recuperação aumentada por geração~\cite{lewis2020rag}. Em \armTextToCypher{} e \armGraphTools{}, o modelo consulta o mesmo grafo e dispõe de no máximo \nMaxCalls{}
chamadas por pergunta, em um ciclo de raciocínio e ação~\cite{yao2023react}. A diferença está em quem escreve a
lógica: no primeiro, o modelo escreve toda consulta~\cite{ozsoy2025text2cypher}, sem a biblioteca de algoritmos de
grafos (GDS)~\cite{neo4j_gds} e sem os padrões táticos já calculados, e a ferramenta só aceita leitura, devolve até
\nMaxRows{} linhas e interrompe a consulta após \nTimeoutS{} s; no segundo, a lógica difícil (jogadas, rede de
passes, intermediação) está em ferramentas escritas e congeladas antes das perguntas que as testam. Esse par
controlado isola o efeito de quem escreve a lógica.

A comparação do par é assimétrica. O \armGraphTools{} foi ajustado em sete versões congeladas, a partir dos erros
observados em rodadas anteriores, e seis perguntas foram reexecutadas depois das últimas correções, de modo que o
resultado reportado é o da segunda execução (antes dela, o braço acertou \preFixKGraphTools{} de
\preFixNGraphTools{} perguntas). O \armTextToCypher{} foi medido em primeiro contato, sem ajuste guiado por erro.
A diferença entre os dois tende, portanto, a superestimar o efeito da lógica pronta.

\subsection{Perguntas, gabarito e métricas}

As \nQuestions{} perguntas, em linguagem de futebol e de resposta fechada (quem, quantos, quais), formam cinco
grupos: busca e contagem (\nGroupSearch{}), relações na rede de passes (\nGroupRelations{}), jogadas e sequências
(\nGroupPlays{}), antes e depois de um momento (\nGroupMoments{}) e controle (\nGroupControl{}). O controle reúne
perguntas cuja resposta não existe nos dados, como as que pedem velocidade; a resposta certa é a ausência de dados.
Os resultados são reportados com e sem ele, isto é, nas \nAnswerable{} perguntas com resposta e nas
\nQuestions{}.

O gabarito foi calculado sem o Neo4j e sem o modelo, com pandas sobre os dados brutos e networkx para as redes, e
a intermediação segue o algoritmo de Brandes~\cite{brandes2001faster}. Cada pergunta foi resolvida em todas as
leituras razoáveis (por exemplo, rede dirigida ou não, passes ou xT como peso). Só permaneceram as perguntas
cuja resposta não muda entre leituras. A conferência compara a resposta do modelo com o gabarito conforme o tipo
(jogador, valor, conjunto, jogador e valor, ou ausência de dados), com nomes normalizados pelos apelidos das
escalações.

Cada execução (uma pergunta em um braço) é classificada como acerto, alucinação (resposta errada dada como
resposta), abstenção ou erro de formato; também se registram tokens, chamadas de ferramenta e latência. O acerto de
um braço é a proporção de perguntas respondidas corretamente, com intervalo de Wilson de \ciLevel{}, tomando a
pergunta como unidade. Os braços respondem às mesmas perguntas, e por isso os pares de degraus vizinhos são
comparados pelo teste exato de McNemar, com ajuste de Holm para os três testes. A rodada tem \nRepeats{} repetição
por pergunta; as conclusões são lidas como indícios, e não como demonstração.

\subsection{Escopo}

Ficam fora do escopo: várias partidas, dados de rastreamento ou fisiológicos, inferência sobre lesão, fadiga ou
saúde, valor de mercado, escalação, predição de resultado, qualidade da explicação tática, comparação entre modelos
e agentes que escrevem código. Os dados são abertos, de uma partida pública, e não incluem dado pessoal sensível;
% TODO autor: confirmar o parágrafo de LGPD (o rascunho do colega não está no repositório)
o trabalho não infere condição de saúde de atleta.
